% Options for packages loaded elsewhere
\PassOptionsToPackage{unicode}{hyperref}
\PassOptionsToPackage{hyphens}{url}
\PassOptionsToPackage{dvipsnames,svgnames,x11names}{xcolor}
%
\documentclass[
  10pt,
  a4paper,
]{article}
\usepackage{amsmath,amssymb}
\usepackage{iftex}
\ifPDFTeX
  \usepackage[T1]{fontenc}
  \usepackage[utf8]{inputenc}
  \usepackage{textcomp} % provide euro and other symbols
\else % if luatex or xetex
  \usepackage{unicode-math} % this also loads fontspec
  \defaultfontfeatures{Scale=MatchLowercase}
  \defaultfontfeatures[\rmfamily]{Ligatures=TeX,Scale=1}
\fi
\usepackage{lmodern}
\ifPDFTeX\else
  % xetex/luatex font selection
\fi
% Use upquote if available, for straight quotes in verbatim environments
\IfFileExists{upquote.sty}{\usepackage{upquote}}{}
\IfFileExists{microtype.sty}{% use microtype if available
  \usepackage[]{microtype}
  \UseMicrotypeSet[protrusion]{basicmath} % disable protrusion for tt fonts
}{}
\makeatletter
\@ifundefined{KOMAClassName}{% if non-KOMA class
  \IfFileExists{parskip.sty}{%
    \usepackage{parskip}
  }{% else
    \setlength{\parindent}{0pt}
    \setlength{\parskip}{6pt plus 2pt minus 1pt}}
}{% if KOMA class
  \KOMAoptions{parskip=half}}
\makeatother
\usepackage{xcolor}
\usepackage[margin=23mm]{geometry}
\usepackage{color}
\usepackage{fancyvrb}
\newcommand{\VerbBar}{|}
\newcommand{\VERB}{\Verb[commandchars=\\\{\}]}
\DefineVerbatimEnvironment{Highlighting}{Verbatim}{commandchars=\\\{\}}
% Add ',fontsize=\small' for more characters per line
\newenvironment{Shaded}{}{}
\newcommand{\AlertTok}[1]{\textcolor[rgb]{1.00,0.00,0.00}{\textbf{#1}}}
\newcommand{\AnnotationTok}[1]{\textcolor[rgb]{0.38,0.63,0.69}{\textbf{\textit{#1}}}}
\newcommand{\AttributeTok}[1]{\textcolor[rgb]{0.49,0.56,0.16}{#1}}
\newcommand{\BaseNTok}[1]{\textcolor[rgb]{0.25,0.63,0.44}{#1}}
\newcommand{\BuiltInTok}[1]{\textcolor[rgb]{0.00,0.50,0.00}{#1}}
\newcommand{\CharTok}[1]{\textcolor[rgb]{0.25,0.44,0.63}{#1}}
\newcommand{\CommentTok}[1]{\textcolor[rgb]{0.38,0.63,0.69}{\textit{#1}}}
\newcommand{\CommentVarTok}[1]{\textcolor[rgb]{0.38,0.63,0.69}{\textbf{\textit{#1}}}}
\newcommand{\ConstantTok}[1]{\textcolor[rgb]{0.53,0.00,0.00}{#1}}
\newcommand{\ControlFlowTok}[1]{\textcolor[rgb]{0.00,0.44,0.13}{\textbf{#1}}}
\newcommand{\DataTypeTok}[1]{\textcolor[rgb]{0.56,0.13,0.00}{#1}}
\newcommand{\DecValTok}[1]{\textcolor[rgb]{0.25,0.63,0.44}{#1}}
\newcommand{\DocumentationTok}[1]{\textcolor[rgb]{0.73,0.13,0.13}{\textit{#1}}}
\newcommand{\ErrorTok}[1]{\textcolor[rgb]{1.00,0.00,0.00}{\textbf{#1}}}
\newcommand{\ExtensionTok}[1]{#1}
\newcommand{\FloatTok}[1]{\textcolor[rgb]{0.25,0.63,0.44}{#1}}
\newcommand{\FunctionTok}[1]{\textcolor[rgb]{0.02,0.16,0.49}{#1}}
\newcommand{\ImportTok}[1]{\textcolor[rgb]{0.00,0.50,0.00}{\textbf{#1}}}
\newcommand{\InformationTok}[1]{\textcolor[rgb]{0.38,0.63,0.69}{\textbf{\textit{#1}}}}
\newcommand{\KeywordTok}[1]{\textcolor[rgb]{0.00,0.44,0.13}{\textbf{#1}}}
\newcommand{\NormalTok}[1]{#1}
\newcommand{\OperatorTok}[1]{\textcolor[rgb]{0.40,0.40,0.40}{#1}}
\newcommand{\OtherTok}[1]{\textcolor[rgb]{0.00,0.44,0.13}{#1}}
\newcommand{\PreprocessorTok}[1]{\textcolor[rgb]{0.74,0.48,0.00}{#1}}
\newcommand{\RegionMarkerTok}[1]{#1}
\newcommand{\SpecialCharTok}[1]{\textcolor[rgb]{0.25,0.44,0.63}{#1}}
\newcommand{\SpecialStringTok}[1]{\textcolor[rgb]{0.73,0.40,0.53}{#1}}
\newcommand{\StringTok}[1]{\textcolor[rgb]{0.25,0.44,0.63}{#1}}
\newcommand{\VariableTok}[1]{\textcolor[rgb]{0.10,0.09,0.49}{#1}}
\newcommand{\VerbatimStringTok}[1]{\textcolor[rgb]{0.25,0.44,0.63}{#1}}
\newcommand{\WarningTok}[1]{\textcolor[rgb]{0.38,0.63,0.69}{\textbf{\textit{#1}}}}
\setlength{\emergencystretch}{3em} % prevent overfull lines
\providecommand{\tightlist}{%
  \setlength{\itemsep}{0pt}\setlength{\parskip}{0pt}}
\setcounter{secnumdepth}{-\maxdimen} % remove section numbering
\ifLuaTeX
  \usepackage{selnolig}  % disable illegal ligatures
\fi
\usepackage{bookmark}
\IfFileExists{xurl.sty}{\usepackage{xurl}}{} % add URL line breaks if available
\urlstyle{same}
\hypersetup{
  colorlinks=true,
  linkcolor={black},
  filecolor={Maroon},
  citecolor={Blue},
  urlcolor={blue},
  pdfcreator={LaTeX via pandoc}}

\author{}
\date{}

\begin{document}

\section{DDTA BA5 canonical referent naming and controlled-registry
candidate -
R1}\label{ddta-ba5-canonical-referent-naming-and-controlled-registry-candidate---r1}

\textbf{Status:} PROVISIONAL CANDIDATE AFTER BA5-T1 / BA5 OPEN

\textbf{Repository baseline reviewed:}
\texttt{aa4b785dc72c2fbcd20fd04976e77fca3d07bf25}

\textbf{Closed dependencies:} documentation layer; BA0 responsibility
boundary; BA1 \texttt{BAReferent\ +\ BAProposition}; BA2 semantic
registries; BA3 provenance/continuity; BA4 projection contract.

\subsection{1. Purpose}\label{purpose}

BA5-T1 replaces the initial synonym-first pressure with a stronger
falsifiable authoring hypothesis:

\begin{quote}
For project referents that are named in governed DDTA material, one
exact canonical name is used across the governed baseline, accepted Base
Analysis and every derived view that presents that referent as shared
project meaning. Synonymous or alternative entity identifiers are not
normative input.
\end{quote}

The motivating example is deliberately simple:

\begin{Shaded}
\begin{Highlighting}[]
\NormalTok{canonical project referent name:}
\NormalTok{  cameraIngresso}
\end{Highlighting}
\end{Shaded}

If \texttt{cameraIngresso} denotes one governed project referent, a
second document or derived shared view does not rename it to
\texttt{entryCamera}, \texttt{cameraEntrata}, \texttt{CameraIngresso} or
another alias.

The purpose is not stylistic uniformity. The purpose is to make
cross-document identity, projection comparison, source drill-down and
tool integration mechanically reliable without requiring synonym
resolution before the Base Analysis architecture itself has been fully
regressed.

\subsection{2. Provisional lower-bound
rule}\label{provisional-lower-bound-rule}

For one governed baseline and naming scope:

\begin{Shaded}
\begin{Highlighting}[]
\NormalTok{one shared named project referent}
\NormalTok{    {-}\textgreater{} exactly one canonicalName token}

\NormalTok{same referent across governed documents}
\NormalTok{    {-}\textgreater{} same exact canonicalName}

\NormalTok{same referent exposed by accepted BA / shared projection rendering}
\NormalTok{    {-}\textgreater{} same exact canonicalName}

\NormalTok{alternative lexical identifier for same referent}
\NormalTok{    {-}\textgreater{} REJECTED as normative authoring}
\end{Highlighting}
\end{Shaded}

Exact means token equality under the declared project convention. Case
or formatting variants do not silently alias:

\begin{Shaded}
\begin{Highlighting}[]
\NormalTok{cameraIngresso}
\NormalTok{CameraIngresso}
\NormalTok{camera\_ingresso}
\NormalTok{entryCamera}
\end{Highlighting}
\end{Shaded}

are distinct lexical tokens. Only the registered canonical token is
accepted for the referent.

\subsection{3. Name is not identity}\label{name-is-not-identity}

The strongest negative control in T1 is a governed rename.

BA1 identity and BA3 continuity already establish that lexical wording
does not define semantic identity. BA3 explicitly permits a
\texttt{BAReferent} to RETAIN when the same independently reusable
methodology-neutral project meaning survives a change in wording.

Therefore BA5 must not make:

\begin{Shaded}
\begin{Highlighting}[]
\NormalTok{canonicalName == BAReferent identity}
\end{Highlighting}
\end{Shaded}

The current lower bound is instead:

\begin{Shaded}
\begin{Highlighting}[]
\NormalTok{BAReferent semantic identity}
\NormalTok{    != canonicalName}

\NormalTok{canonicalName}
\NormalTok{    is baseline{-}scoped governed lexical metadata}
\NormalTok{    used consistently wherever that referent is named}
\end{Highlighting}
\end{Shaded}

Example:

\begin{Shaded}
\begin{Highlighting}[]
\NormalTok{B0:}
\NormalTok{  same semantic referent R{-}camera{-}entry}
\NormalTok{  canonicalName = cameraIngresso}

\NormalTok{B1 governed rename:}
\NormalTok{  same semantic referent retained by BA3}
\NormalTok{  canonicalName = cameraNord}
\end{Highlighting}
\end{Shaded}

Historical B0 material remains \texttt{cameraIngresso}; current B1
governed documents, BA rendering and rebuilt shared projections use
\texttt{cameraNord}.

A rename is therefore a governed lexical change, not automatically a new
BAE identity.

\subsection{4. Canonical-name registry
responsibility}\label{canonical-name-registry-responsibility}

BA5-T1 requires a governed canonical-name registry responsibility, but
does not mandate one physical database, JSON registry, graph node family
or ThreatForge table.

The semantic minimum is:

\begin{Shaded}
\begin{Highlighting}[]
\NormalTok{CanonicalReferentNameRegistry             [governed lexical responsibility; NOT BAE]}
\NormalTok{|{-} governedBaselineKey               1}
\NormalTok{|{-} namingScopeKey                    1}
\NormalTok{\textasciigrave{}{-} entry                             0..*}
\NormalTok{     |{-} canonicalName                1 exact token}
\NormalTok{     \textasciigrave{}{-} governedReferentBinding      1 unambiguous within baseline/scope}
\end{Highlighting}
\end{Shaded}

An accepted BA materialization must be able to resolve a named
\texttt{BAReferent} to that canonical token without inventing a synonym.

The registry may be realized as governed authoring metadata, a generated
validated symbol table over governed files, or another auditable
mechanism. The physical realization is deferred.

\subsection{5. Registry authority
boundary}\label{registry-authority-boundary}

The registry does not make Base Analysis or ThreatForge project
authority.

The authority chain remains:

\begin{Shaded}
\begin{Highlighting}[]
\NormalTok{governed project authoring / canonical{-}name registration}
\NormalTok{    {-}\textgreater{} governed documentation baseline}
\NormalTok{        {-}\textgreater{} accepted BA materialization}
\NormalTok{            {-}\textgreater{} derived views}
\end{Highlighting}
\end{Shaded}

A tool may enforce or validate the registered name, but may not silently
decide that an unregistered alias is equivalent.

If the author attempts:

\begin{Shaded}
\begin{Highlighting}[]
\NormalTok{entryCamera}
\end{Highlighting}
\end{Shaded}

when the registered referent is:

\begin{Shaded}
\begin{Highlighting}[]
\NormalTok{cameraIngresso}
\end{Highlighting}
\end{Shaded}

the minimum valid behavior is rejection or an explicit correction
proposal. Automatic hidden alias normalization is outside the accepted
T1 core.

\subsection{6. Cross-document naming
invariant}\label{cross-document-naming-invariant}

For a named referent, every semantic reference in governed DDTA
documents uses the canonical token.

Allowed:

\begin{Shaded}
\begin{Highlighting}[]
\NormalTok{The component \textasciigrave{}cameraIngresso\textasciigrave{} transfers ...}
\NormalTok{Decision D{-}X constrains \textasciigrave{}cameraIngresso\textasciigrave{} ...}
\NormalTok{FR{-}Y requires \textasciigrave{}cameraIngresso\textasciigrave{} ...}
\end{Highlighting}
\end{Shaded}

Rejected as an alternative identifier for the same referent:

\begin{Shaded}
\begin{Highlighting}[]
\NormalTok{entry camera}
\NormalTok{cameraEntrata}
\NormalTok{CameraIngresso}
\end{Highlighting}
\end{Shaded}

Free explanatory prose may surround the canonical token, but an
alternative phrase must not become a second machine-significant
identifier for the same referent.

BA5-T1 does not attempt automatic migration of arbitrary legacy
narrative prose.

\subsection{7. Uniqueness lower bound}\label{uniqueness-lower-bound}

Two different shared referents in the same governed baseline/naming
scope cannot resolve to the same canonical token.

\begin{Shaded}
\begin{Highlighting}[]
\NormalTok{referent A {-}\textgreater{} cameraIngresso}
\NormalTok{referent B {-}\textgreater{} cameraIngresso}
\end{Highlighting}
\end{Shaded}

is rejected because consumers could not know which identity a source or
projection reference denotes.

The lower bound therefore requires canonical-name uniqueness within a
declared naming scope.

T1 does not yet freeze whether a future implementation uses one
project-global namespace or explicit sub-namespaces. That is a
representation/governance detail for later pressure if concrete projects
require it.

\subsection{8. Shared projection
rendering}\label{shared-projection-rendering}

BA4 permits presentation labels that preserve shared meaning. BA5-T1
narrows that freedom for shared referent naming:

\begin{quote}
A projection may choose its own layout, local item kind, grouping and
method vocabulary, but when it displays a traced BAReferent as the
shared project entity, it uses the referent's canonical name.
\end{quote}

Allowed:

\begin{Shaded}
\begin{Highlighting}[]
\NormalTok{[FlowParticipant] cameraIngresso}
\NormalTok{[AssuranceSubject] cameraIngresso}
\end{Highlighting}
\end{Shaded}

Here \texttt{FlowParticipant} and \texttt{AssuranceSubject} are
projection-owned kinds; \texttt{cameraIngresso} remains the shared
entity name.

Rejected when presented as the same shared entity:

\begin{Shaded}
\begin{Highlighting}[]
\NormalTok{[FlowParticipant] Entry Camera}
\NormalTok{[AssuranceSubject] Ingress Sensor}
\end{Highlighting}
\end{Shaded}

The method may still create a genuinely method-owned interpretation item
such as:

\begin{Shaded}
\begin{Highlighting}[]
\NormalTok{external{-}ingress{-}exposure}
\end{Highlighting}
\end{Shaded}

provided that item is explicitly method-owned, trace-bound and not
presented as a replacement name for \texttt{cameraIngresso}.

\subsection{9. Aggregation and local method
items}\label{aggregation-and-local-method-items}

A projection-local item need not correspond one-to-one with a
BAReferent. A method may aggregate several BA elements under a local
analysis item.

In that case the local item may have a method-owned label because it is
not claiming to be the canonical project entity.

The boundary is:

\begin{Shaded}
\begin{Highlighting}[]
\NormalTok{label denotes shared BAReferent}
\NormalTok{    {-}\textgreater{} canonicalName required}

\NormalTok{label denotes projection{-}owned aggregation/interpretation}
\NormalTok{    {-}\textgreater{} local label allowed}
\NormalTok{       + BA trace required}
\NormalTok{       + must not masquerade as shared referent identity}
\end{Highlighting}
\end{Shaded}

This preserves BA4 method freedom without sacrificing shared entity
naming consistency.

\subsection{10. Operator, role and semantic-kind
keys}\label{operator-role-and-semantic-kind-keys}

BA2 already closes stable canonical semantic keys such as:

\begin{Shaded}
\begin{Highlighting}[]
\NormalTok{transfer}
\NormalTok{consumeService}
\NormalTok{assignResponsibility}
\NormalTok{source}
\NormalTok{destination}
\NormalTok{content}
\NormalTok{store}
\NormalTok{contract}
\end{Highlighting}
\end{Shaded}

The same controlled-registry philosophy is compatible with those keys,
but BA5-T1 does not reopen or redesign their registries.

T1 is intentionally limited to the referent-name boundary. The next BA5
pressure target may test whether the same exact-token/no-synonym
authoring rule can cover the complete operative BA vocabulary and
registry-extension workflow without semantic loss.

\subsection{11. Tool-support boundary}\label{tool-support-boundary}

The initial tool hypothesis is intentionally narrow.

A tool may:

\begin{itemize}
\tightlist
\item
  validate exact canonical-name use;
\item
  prevent duplicate names in the active naming scope;
\item
  offer completion from the accepted registry;
\item
  flag an unregistered token;
\item
  present a candidate correction or candidate registry-extension
  request.
\end{itemize}

A tool may not, by default:

\begin{itemize}
\tightlist
\item
  create a hidden synonym dictionary;
\item
  silently map an alias to a canonical referent;
\item
  mint a new accepted referent name without governance;
\item
  use method taxonomy as a replacement project name; or
\item
  make a lexical guess into project or BA truth.
\end{itemize}

This keeps ThreatForge or any future tool in an enforcement/assistance
role rather than semantic authority.

\subsection{12. Mutation pressure}\label{mutation-pressure}

\subsubsection{N1 - second document introduces an
alias}\label{n1---second-document-introduces-an-alias}

\begin{Shaded}
\begin{Highlighting}[]
\NormalTok{D1: cameraIngresso}
\NormalTok{D2: entryCamera}
\end{Highlighting}
\end{Shaded}

If both intend the same referent, \texttt{entryCamera} is rejected as
normative authoring. The author must use \texttt{cameraIngresso} or
perform a governed rename.

Result: \textbf{PASS}.

\subsubsection{N2 - case/format
variation}\label{n2---caseformat-variation}

\begin{Shaded}
\begin{Highlighting}[]
\NormalTok{cameraIngresso}
\NormalTok{CameraIngresso}
\NormalTok{camera\_ingresso}
\end{Highlighting}
\end{Shaded}

Only the registered exact token is accepted.

Result: \textbf{PASS}.

\subsubsection{N3 - derived human view renames the
entity}\label{n3---derived-human-view-renames-the-entity}

\begin{Shaded}
\begin{Highlighting}[]
\NormalTok{BAReferent {-}\textgreater{} cameraIngresso}
\NormalTok{view label {-}\textgreater{} Entry Camera}
\end{Highlighting}
\end{Shaded}

If the label claims to name the shared entity, reject it. A descriptive
section heading may differ, but the entity reference itself remains
canonical.

Result: \textbf{PASS}.

\subsubsection{N4 - method projection uses local type
vocabulary}\label{n4---method-projection-uses-local-type-vocabulary}

\begin{Shaded}
\begin{Highlighting}[]
\NormalTok{FlowParticipant(cameraIngresso)}
\NormalTok{AssuranceSubject(cameraIngresso)}
\end{Highlighting}
\end{Shaded}

The local kind may vary while the entity name remains exact.

Result: \textbf{PASS}.

\subsubsection{N5 - governed rename across
baselines}\label{n5---governed-rename-across-baselines}

\begin{Shaded}
\begin{Highlighting}[]
\NormalTok{B0 cameraIngresso}
\NormalTok{B1 cameraNord}
\end{Highlighting}
\end{Shaded}

If the same project meaning survives, BA3 may RETAIN referent identity
while B1 adopts the new canonical name. Historical B0 remains unchanged.

Result: \textbf{PASS WITH REFINEMENT} - canonical name is
baseline-scoped and must not be used as BAE identity.

\subsubsection{N6 - two referents request the same
token}\label{n6---two-referents-request-the-same-token}

Within one naming scope, reject the collision or require distinct
canonical names before the governed baseline is accepted.

Result: \textbf{PASS}.

\subsection{13. Reopen checks}\label{reopen-checks}

BA5-T1 finds no material reason to reopen closed BA responsibilities.

\begin{itemize}
\tightlist
\item
  BA1 remains sufficient: canonical naming is metadata over
  \texttt{BAReferent}, not a third identity family.
\item
  BA2 remains sufficient: naming does not require a \texttt{nameOf}
  operator.
\item
  BA3 already handles retained identity under wording/name change across
  baselines.
\item
  BA4 remains sufficient: the new rule narrows permissible shared
  rendering labels but does not alter projection identity, coverage,
  trace or interpretation mechanics.
\item
  The closed document metamodel need not be reopened: this is a
  cross-document controlled-authoring/governance constraint, not a new
  MR/Decision/FR/SR structural field family.
\end{itemize}

\subsection{14. Provisional
dispositions}\label{provisional-dispositions}

\begin{Shaded}
\begin{Highlighting}[]
\NormalTok{one named referent {-}\textgreater{} one canonicalName per baseline/scope       REQUIRED}
\NormalTok{exact canonical token in governed semantic references           REQUIRED}
\NormalTok{exact canonical token in shared derived referent rendering      REQUIRED}
\NormalTok{synonymous entity identifiers as normative authoring             REJECTED}
\NormalTok{case/format variants as implicit aliases                         REJECTED}
\NormalTok{canonical{-}name collision within naming scope                     REJECTED}
\NormalTok{canonicalName == BAReferent identity                              REJECTED}
\NormalTok{governed rename + BA3 RETAIN                                     PASS}
\NormalTok{projection{-}owned type/category labels                            ALLOWED}
\NormalTok{method{-}owned non{-}referent item labels                            ALLOWED DOWNSTREAM}
\NormalTok{tool exact validation/completion                                 ALLOWED}
\NormalTok{tool hidden synonym normalization                                REJECTED IN T1 CORE}
\NormalTok{new BAE family                                                   NOT FORCED}
\NormalTok{new BA2 operator                                                 NOT FORCED}
\NormalTok{BA1 / BA2 / BA3 / BA4 reopen                                    NOT TRIGGERED}
\NormalTok{documentation metamodel reopen                                   NOT TRIGGERED}
\NormalTok{BA5                                                              STARTED / NOT CLOSED}
\end{Highlighting}
\end{Shaded}

\subsection{15. Falsification rules}\label{falsification-rules}

Revise this candidate if concrete controlled-authoring/integration
evidence demonstrates that:

\begin{enumerate}
\def\labelenumi{\arabic{enumi}.}
\tightlist
\item
  one shared referent must have two concurrent normative names in the
  same baseline/scope to preserve necessary project meaning;
\item
  exact canonical naming makes a required methodology projection
  impossible rather than merely less convenient;
\item
  a canonical rename cannot be handled through existing BA3 continuity
  without semantic identity loss;
\item
  the registry cannot enforce uniqueness without introducing a missing
  shared identity responsibility;
\item
  method-specific labels must replace canonical entity names to preserve
  method semantics; or
\item
  a materially different corpus cannot be authored/integrated without
  synonym equivalence in the semantic core.
\end{enumerate}

Until such evidence exists, synonym machinery remains unnecessary
complexity.

\subsection{16. Smallest unresolved set after
BA5-T1}\label{smallest-unresolved-set-after-ba5-t1}

The next smallest BA5 question is no longer general synonym handling. It
is registry coverage and extension:

\begin{enumerate}
\def\labelenumi{\arabic{enumi}.}
\tightlist
\item
  test exact canonical authoring across operators, roles, semantic kinds
  and controlled values already closed by BA2;
\item
  define the minimum governed extension workflow when a genuinely new
  canonical term is required;
\item
  verify tool validation/completion can enforce the registry without
  becoming semantic authority; and
\item
  determine whether any concrete integration case forces
  aliases/synonyms despite the no-synonym default.
\end{enumerate}

Do not start BA6 until BA5 is explicitly closed.

\end{document}
